\documentclass[11pt,a4paper]{article}

\usepackage{fontspec}
\usepackage[main=russian, english]{babel}
\usepackage{amsmath, amsfonts, amssymb}
\usepackage{graphicx}
\usepackage{booktabs}
\usepackage{cite}
\usepackage{geometry}
\usepackage{xcolor}
\usepackage{url}
\usepackage{hyperref}

% Настройка шрифтов для XeLaTeX/LuaLaTeX (по умолчанию используются системные)
\setmainfont{Times New Roman}
\setromanfont{Times New Roman}
\setsansfont{Arial}
\setmonofont{Courier New}

\geometry{margin=2.5cm}

\title{\textbf{Адаптивное извлечение секвенограмм: совместная слепая деконволюция с помощью физически обоснованной параметрической ConvNMF и Евклидовой тригонометрии потерь}}
\author{Ваше Имя\thanks{Автор для корреспонденции: email@example.com} \\ 
        \small \textit{Кафедра вычислительной обработки сигналов}}
\date{\today}

\begin{document}

\maketitle

\begin{abstract}
Анализ акустических паттернов в естественных средах осложнен многолучевым распространением и перекрытием сигналов. Традиционные методы деконволюции страдают от проблемы «курицы и яйца» при попытке разделить неизвестный источник и неизвестное эхо. В данной работе представлена физически обоснованная генеративная архитектура для извлечения \textit{секвенограмм} — разреженных векторов временной активации индивидуальных звуковых событий. Предложенный алгоритм выполняет совместную (joint) оптимизацию трех компонентов: причинного 1D-профиля эха, инстанс-специфичных секвенограмм и параметрического словаря шаблонов. Для моделирования словаря мы вводим дифференцируемый физический рендерер, где параметры частотной модуляции и гармоник (обертонов) задаются проекциями в касательном пространстве сферы ($T_p S^n$). Для обеспечения сходимости этой сложной многозадачной системы без использования эвристических гиперпараметров, мы представляем концепцию \textit{Локальной калибровочной нормировки} и \textit{Евклидовой тригонометрии градиентов}, которая динамически маршрутизирует вычислительные векторы оптимизатора. Метод демонстрирует исключительную точность при разделении спектрально идентичных, но сдвинутых по времени сигналов в условиях сильной реверберации.
\end{abstract}

\section{Введение}
Распознавание сигналов на частотно-временных представлениях (спектрограммах) в естественных условиях (леса, пещеры) сталкивается с фундаментальной физической проблемой: за первичной эмиссией следует реверберационный «хвост», размывающий акустическую энергию. При высокой плотности сигналов (например, эхолокация стаи летучих мышей) хвосты эха и последующие излучения сливаются в непрерывный спектральный шум.

Стандартные подходы (сверточные нейронные сети или классическая NMF) пытаются решить эту проблему, используя жесткие 2D-шаблоны. Однако такие методы обладают серьезными недостатками. Во-первых, они не способны разделить идентичные перекрывающиеся сигналы (instance segmentation). Во-вторых, обучение словаря паттернов на реальных данных неизбежно приводит к впитыванию реверберации в сам шаблон.

Для решения этих проблем мы предлагаем архитектуру параметрической слепой деконволюции. Метод одновременно извлекает адаптивное эхо среды, локализует звуковые объекты во времени (секвенограммы) и восстанавливает их безэховую спектральную структуру, используя физические законы распространения и генерации звука.

\section{Материалы и методы}

\subsection{Генеративная модель совместной оптимизации}
Пусть $S \in \mathbb{R}_{\ge 0}^{F \times T}$ — наблюдаемая спектрограмма. Мы моделируем ее как результат двух последовательных физических процессов: излучения и реверберации.
Вектор-матрица активаций $H \in \mathbb{R}_{\ge 0}^{K \times T}$ (секвенограмма) описывает точное время звучания $K$ независимых каналов. «Чистая» спектрограмма $V$ формируется сверткой секвенограмм с адаптивными шаблонами $W_k$:
\begin{equation}
V(f, t) = \sum_{k=1}^{K} (H_k *_{t} W_k(f, \cdot))(t)
\end{equation}
Затем $V$ подвергается воздействию причинного частотно-зависимого 1D-фильтра эха $E \in \mathbb{R}_{\ge 0}^{F \times L_e}$:
\begin{equation}
\hat{S}(f, t) = (V(f, \cdot) *_{t} E(f, \cdot))(t)
\end{equation}
Изоляция осей (свертка эха происходит строго вдоль оси времени) математически предотвращает нефизичный перенос энергии между частотами.

Поскольку $H$, $W$ и $E$ неизвестны, их оптимизация должна происходить совместно. В качестве инициализации алгоритма $H$ задается как локальная огибающая энергии (RMS), $E$ инициализируется как экспоненциальное затухание, а $W$ генерируется дифференцируемым рендерером.

\subsection{Физически обоснованный словарь: Дифференцируемый рендеринг и Касательное пространство}
Вместо обучения «сырых» пикселей матрицы $W_k$, мы моделируем словарь как параметрическую физическую кривую. Базовая форма $k$-го канала задается макро-якорем на гиперсфере $\vec{n}_k \in S^n$. Для моделирования микро-флуктуаций (естественной вариативности формы сигнала во времени), нейросеть предсказывает вектор отклонений $\vec{v}_k(t)$, который проецируется на касательную плоскость:
\begin{equation}
\vec{\Delta}_k(t) = \vec{v}_k(t) - (\vec{v}_k(t) \cdot \vec{n}_k)\vec{n}_k
\end{equation}
Экспоненциальное отображение возвращает точку обратно на сферу: $\vec{n}_k(t) = Retract(\vec{n}_k, \vec{\Delta}_k(t))$.

Вектор $\vec{n}_k(t)$ декодируется в параметры физического сплайна основной частоты $f_0(t)$ и независимые амплитудные огибающие $A_1(t), A_2(t)$. Итоговый 2D-шаблон $W_k$ генерируется с помощью оценки плотности (Gaussian Splatting):
\begin{equation}
W_k(f, \tau) = A_1(\tau) \exp\left( - \frac{(f - f_0(\tau))^2}{2\sigma^2} \right) + A_2(\tau) \exp\left( - \frac{(f - 2f_0(\tau))^2}{2\sigma^2} \right)
\end{equation}
Данный подход обладает тремя фундаментальными преимуществами: 1) исключает появление артефактов (в отличие от сферических гармоник); 2) гарантирует жесткую физическую привязку первого обертона ($2f_0$); 3) позволяет оптимизировать гладкое латентное многообразие кривых.

\subsection{Инстанс-сегментация: Локальная нормировка и Евклидова тригонометрия}
Чтобы слепая деконволюция корректно разделяла перекрывающиеся сигналы, необходимо наложить структурные ограничения. Для извлечения секвенограмм мы формулируем три физических условия: бинаризацию ($\mathcal{L}_{sharp}$), временную локализацию (штраф дисперсии $\mathcal{L}_{loc}$, заставляющий каждый канал описывать строго один физический объект) и разреженность каналов ($\mathcal{L}_{sparse}$).

Традиционное многозадачное обучение объединяет штрафы линейно ($\sum \alpha_i \mathcal{L}_i$), что требует ручного подбора гиперпараметров и часто приводит к затуханию градиентов в разреженных сигналах. Мы предлагаем агрегировать компоненты с использованием \textbf{Евклидовой метрики} $\mathcal{L}_{reg} = \sqrt{\mathcal{L}_{sharp}^2 + \mathcal{L}_{loc}^2 + \mathcal{L}_{sparse}^2}$.

Чтобы Евклидова норма была математически корректной, мы применяем принцип \textbf{Локальной калибровочной нормировки}. Мы отказываемся от глобального усреднения градиентов. Калибровка производится \textit{интенсивно}: внутри зоны локального взаимодействия (вблизи $h(t) > 0$) амплитуда вариационной производной каждого функционала аналитически приводится к единице: $\left| \frac{\delta \mathcal{L}}{\delta h} \right|_{local} = 1$.

Это формирует изотропное пространство ошибок. Рассмотрим сигнал градиента для двумерного случая $\mathcal{L}_{reg} = \sqrt{\mathcal{L}_A^2 + \mathcal{L}_B^2}$. По правилу дифференцирования сложной функции:
\begin{equation}
\frac{\partial \mathcal{L}_{reg}}{\partial h_i} = \frac{\mathcal{L}_A}{\sqrt{\mathcal{L}_A^2 + \mathcal{L}_B^2}} \frac{\partial \mathcal{L}_A}{\partial h_i} + \frac{\mathcal{L}_B}{\sqrt{\mathcal{L}_A^2 + \mathcal{L}_B^2}} \frac{\partial \mathcal{L}_B}{\partial h_i}
\end{equation}
Введя макроскопический угол состояния системы $\theta = \arctan (\mathcal{L}_B / \mathcal{L}_A)$, получаем уравнение градиентного спуска в тригонометрической форме:
\begin{equation}
\frac{\partial \mathcal{L}_{reg}}{\partial h_i} = \underbrace{\cos(\theta)}_{\text{глобальный вес}} \cdot \underbrace{\frac{\partial \mathcal{L}_A}{\partial h_i}}_{\text{откалиброванная сила}} + \underbrace{\sin(\theta)}_{\text{глобальный вес}} \cdot \underbrace{\frac{\partial \mathcal{L}_B}{\partial h_i}}_{\text{откалиброванная сила}}
\end{equation}

Это уравнение раскрывает элегантную двухуровневую систему маршрутизации градиентов:
\begin{itemize}
    \item \textbf{Микроскопический баланс:} В точках конфликта амплитуды локальных сил равны $|1|$. Фазовое пространство не искажается размерностями штрафов.
    \item \textbf{Макроскопический авто-баланс:} Вектор приоритетов не фиксирован гиперпараметрами, а динамически вращается. Если сеть минимизировала проблему $A$ ($\cos(\theta) \to 0$), оптимизатор автоматически переключает $100\%$ вычислительного внимания ($\sin(\theta) \to 1$) на проблему $B$.
\end{itemize}

\section{Результаты}

\subsection{Синтетическая валидация (Разделение инстансов)}
Для проверки концепции совместной оптимизации (Joint Optimization) был сгенерирован набор данных, содержащий перекрывающиеся спектрально идентичные вызовы (имитация стаи) с сильным частотно-зависимым эхом.
Модель успешно избежала тривиальных решений. Благодаря штрафу дисперсии ($\mathcal{L}_{loc}$), идентичные, но разнесенные во времени сигналы были классифицированы в разные каналы $H_k$. Дифференцируемый рендерер точно восстановил независимые огибающие амплитуд основного тона и первого обертона, а модель эха $E$ автоматически выявила высокие скорости затухания для высоких частот.

\section{Обсуждение}
Главным достижением предложенной архитектуры является физически обоснованное устранение проблемы «курицы и яйца» в задачах деконволюции. Жесткая параметризация библиотеки (сплайны и гармоники $2f_0$) в сочетании со штрафом микро-флуктуаций в касательном пространстве $T_p S^n$ не позволяет модели впитывать реверберацию в шаблон. Оставшаяся дисперсия акустической энергии математически вынуждена абсорбироваться 1D-матрицей эха $E$.

Кроме того, мы представили концепцию Евклидовой тригонометрии градиентов. Обеспечив единичную вариационную силу на локальном (микроскопическом) уровне, мы получили систему, которая самостоятельно вычисляет макроскопический угол атаки $\theta$ на каждом шаге оптимизации, полностью устраняя необходимость ручной балансировки многозадачных функций потерь.

\section{Заключение}
Представленный фреймворк параметрической ConvNMF предлагает строгий физический подход к анализу реверберирующих спектрограмм. Сочетание дифференцируемого гауссова рендеринга, проекций на сферу и масштабно-инвариантной евклидовой регуляризации превращает метод в вычислительно эффективный и полностью автономный инструмент для обнаружения паттернов, инстанс-сегментации и анализа среды в задачах крупномасштабного акустического мониторинга.

\section*{Доступность данных и кода}
Исходный код (PyTorch), включая реализацию FFT-сверток, дифференцируемого рендеринга и калиброванных функций потерь, доступен в нашем репозитории: \url{https://github.com/yourusername/AdaptiveSequenograms}.

\bibliographystyle{plain}
\bibliography{references}

\end{document}